\begin{frame}{두 번째 실수: ``평균 에피소드 길이''로 방문 횟수
추정하기}
  위 표의 마지막 열 --- $(0,0)$ 은 50,546회, $(0,2)$ 는
  100,155회 방문됨. ``에피소드당 평균 길이가 대략 5.3스텝이고
  20만 에피소드이므로 모든 상태가 대략 같은 횟수($\approx$10만)
  를 방문했을 것''이라고 \textbf{생각하기} 쉽다.
  \vskip0.6em
  그런데 실제로는 \textbf{어떤 상태가 ``몇 에피소드에서 몇
  번'' 방문되느냐}가 상태마다 다르다 --- $(0,2)$ 는 ``위쪽
  루프를 타는 에피소드''마다 2번 이상(시작 + 루프 중간)
  나타나고, $(0,0)$ 은 ``아래쪽 루프 에피소드''에서는 0번일 수
  있다.
  \vskip0.6em
  \begin{block}{원칙: 방문 횟수 = 신뢰도}
    (5.2절 블랙잭에서 보았던 그 원칙) 이 ``실제 횟수''는
    \textbf{계산해봐야} 알 수 있다. ``평균 길이 $\times$ 에피소드
    수''로 \textbf{대충 세는} 것은 방문이 \textbf{불균형한}
    환경(루프가 여러 개인 잔디깎는 기, 낮은 합이 드문 블랙잭)에서
    신뢰도를 \textbf{잘못 판단}하게 만든다.
  \end{block}
\end{frame}
